\section{SF-E814: an actual simple degree-two cone at the sharp secant
Euler
value}\label{SF-E814}

Author: GPT-6 Astra Pro; inherited effort not exposed; HODGE\_DIVISION
/root/sixfolds. Date: 2026-11-21.

\textbf{Status:} Whole-proof mathematical PASS received from /root on
2026-11-21 (GPT-6 Astra Pro; inherited effort not exposed), after
independent reconstruction of all three long exact sequences and full
trace consequences. This constructs a genuine ordinary derived object of
rational two-pure secant type. It is not a balanced-RM seed and does not
close the split-eightfold gate.

\subsection{Statement and
construction}\label{SF-E814-statement-and-construction}

Let \(X\) be a complex abelian fourfold, and let \(A,B\) be line bundles
for which

\[
 \delta=c_1(B)-c_1(A)
\]

has a nondegenerate Hermitian form of signature \((2,2)\). Put

\[
 m=\chi(A^\vee\otimes B)=\int_X\frac{\delta^4}{24}>0.
\]

The published Mumford vanishing/index theorem, with the exact locators
recorded in SF-E813, gives

\[
 \operatorname{Ext}^j(A,B)=\operatorname{Ext}^j(B,A)=0\quad(j\ne2),
 \qquad \dim\operatorname{Ext}^2(A,B)=m.
 \tag{E814.1}
\]

Choose any nonzero \(e\in\operatorname{Ext}^2(A,B)\), and define \(F\)
by the distinguished triangle

\[
 B[1]\xrightarrow{i}F\xrightarrow{p}A\xrightarrow{e}B[2].
 \tag{E814.2}
\]

This is an actual perfect complex. Its only nonzero cohomology sheaves
are \(\mathcal H^{-1}(F)=B\) and \(\mathcal H^0(F)=A\), and its
Postnikov invariant is \(e\). It has

\[
 \operatorname{ch}(F)=e^{c_1(A)}-e^{c_1(B)}.
 \tag{E814.3}
\]

Then \(F\) is simple and gluable, and its self-Ext dimensions are

\[
 (\operatorname{ext}^0,\operatorname{ext}^1,
   \operatorname{ext}^2,\operatorname{ext}^3,
   \operatorname{ext}^4)
   =(1,m+7,12,m+7,1).
 \tag{E814.4}
\]

Its full degree-two trace is injective. Its full degree-one trace has
rank eight and kernel dimension \(m-1\). Thus when \(m=1\), the full
trace is injective in degrees zero, one and two. In that case
\(F\boxtimes F^\vee\) on \(X\times X\) is also simple, gluable and fully
semiregular in degree two.

\subsection{1. Direct calculation of the lower self-Ext
groups}\label{SF-E814-direct-calculation-of-the-lower-self-ext-groups}

For a line bundle on an abelian fourfold,

\[
 \operatorname{Ext}^j(A,A)=\operatorname{Ext}^j(B,B)
   =H^j(X,\mathcal O_X),\qquad h^j(X,\mathcal O_X)=\binom4j.
 \tag{E814.5}
\]

Apply \(\operatorname{RHom}(A,-)\) to (E814.2). The only cross group is
\(\operatorname{Ext}^2(A,B)\). The connecting map

\[
 \operatorname{Hom}(A,A)\longrightarrow\operatorname{Ext}^2(A,B),
 \qquad 1\longmapsto e
\]

is injective. The resulting long exact sequence gives

\[
 \begin{array}{c|ccccc}
 j&0&1&2&3&4\\ \hline
 \dim\operatorname{Ext}^j(A,F)&0&m+3&6&4&1.
 \end{array}
 \tag{E814.6}
\]

There are no negative groups. In particular the degree-one group fits
into the exact sequence with left term
\(\operatorname{Ext}^2(A,B)/\mathbf Ce\), of dimension \(m-1\), and
right term \(H^1(\mathcal O_X)\), of dimension four.

Apply \(\operatorname{RHom}(B[1],-)\) to the same triangle. Since
\(\operatorname{RHom}(B[1],A)\) has its sole cohomology in degree three,
the maps induced by \(i\) give isomorphisms

\[
 \operatorname{Ext}^j(B,B)
       \xrightarrow{\sim}\operatorname{Ext}^j(B[1],F)
       \quad(j=0,1,2).
 \tag{E814.7}
\]

Again there are no negative groups.

Finally apply \(\operatorname{RHom}(-,F)\). The relevant connecting maps
are

\[
 \partial_j:\operatorname{Ext}^j(B[1],F)
              \longrightarrow\operatorname{Ext}^{j+1}(A,F),
 \tag{E814.8}
\]

obtained by precomposing with the boundary \(e\), with the indicated
shifts. All three maps \(\partial_0,\partial_1,\partial_2\) are zero.
For \(j=0\), the generator of (E814.7) is \(i\), whose composite with
the preceding boundary vanishes by the distinguished-triangle identity.
For \(j=1,2\), (E814.7) represents every class as \(i\) composed with a
class \(a\in\operatorname{Ext}^j(B,B)\). The corresponding connecting
class factors through \(a\circ e\in\operatorname{Ext}^{j+2}(A,B)=0\), by
(E814.1).

Thus the long exact sequence gives

\[
 \operatorname{Hom}(F,F)=\mathbf C,
\]

and short exact sequences with dimensions

\[
 \operatorname{ext}^1(F,F)=(m+3)+4=m+7,
 \qquad
 \operatorname{ext}^2(F,F)=6+6=12.
 \tag{E814.9}
\]

Both outer complexes in the triangle for \(\operatorname{RHom}(-,F)\)
have no negative cohomology, so neither does
\(\operatorname{RHom}(F,F)\). This proves gluability. Serre duality on
the fourfold then supplies all remaining groups in (E814.4), and gives
zero above degree four.

This calculation uses actual distinguished triangles and long exact
sequences. It does not assume formality of the full dg or \(A_\infty\)
category generated by the two line bundles.

\subsection{2. The complete trace
calculation}\label{SF-E814-the-complete-trace-calculation}

The Fock four-vectors of the two exponential terms in (E814.3) are
decomposable. Their supporting four-planes are graphs whose difference
is the Hermitian form of \(\delta\); nondegeneracy of \(\delta\) makes
them transverse. Therefore the complete two-pure secant calculation of
SF-E801 applies, with two nonzero coefficients. In particular

\[
 \operatorname{rank}c^1_{\operatorname{ch}(F)}=8,
 \qquad
 \operatorname{rank}c^2_{\operatorname{ch}(F)}=12.
 \tag{E814.10}
\]

Evaluation followed by full trace is cap product. Its degree-two rank is
twelve, while \(\operatorname{Ext}^2(F,F)\) has dimension twelve by
(E814.9). Consequently the full degree-two trace is injective, and the
degree-two evaluation is surjective.

For completeness, the degree-one trace rank is not being inferred from
the rank of evaluation alone. The argument of SF-E810 gives the needed
upper bound. Let \(L,M\) be the transverse supporting four-planes of the
two pure terms. The degree-two cap kernel is \(I=L\otimes M\). Because
the degree-two trace is injective, this is also the kernel of degree-two
evaluation. The full trace/cap module identity gives, for every
\(z\in\operatorname{Ext}^1(F,F)\) and \(\xi\in I\),

\[
 \xi\wedge\sigma^1(z)=\sigma^3(\rho^2(\xi)\circ z)=0.
\]

The common annihilator of \(I\) in Fock degree five is

\[
 (\bigwedge^4L\otimes M)\oplus(L\otimes\bigwedge^4M),
\]

of dimension eight, by the explicit exterior calculation in SF-E810.
Thus \(\operatorname{rank}\sigma^1\le8\); (E814.10) gives the reverse
inequality. Hence

\[
 \operatorname{rank}\sigma^1=8,
 \qquad \dim\ker\sigma^1=m-1.
 \tag{E814.11}
\]

The degree-zero trace of the identity is the nonzero character (E814.3),
and \(F\) is simple, so degree zero is injective as well. Finally

\[
 \chi(F,F)=2-2(m+7)+12=-2m,
\]

agreeing directly with Riemann--Roch applied to (E814.3). The case
\(m=1\) realizes exactly the necessary value \(\chi=-2\) in SF-E810; it
shows that the exceptional numerical value in that obstruction is
attained by an actual derived object.

\subsection{3. An explicit CM example and its external
product}\label{SF-E814-an-explicit-cm-example-and-its-external-product}

Take \(X=E_{\mathrm i}^{\,4}\), with its product principal polarization.
On each elliptic curve let \(\mathcal O(0)\) be the degree-one line
bundle of its origin, and put

\[
 A=\mathcal O_X,\qquad
 B=\mathcal O(0)\boxtimes\mathcal O(0)
       \boxtimes\mathcal O(-0)\boxtimes\mathcal O(-0).
 \tag{E814.12}
\]

The associated Hermitian matrix is \(S=\operatorname{diag}(1,1,-1,-1)\).
Its signature is \((2,2)\), and \(\chi(B)=\det S=1\). One can see the
unique nonzero cross Ext group without an additional theorem: the two
positive degree-one elliptic factors have only one-dimensional \(H^0\),
the negative factors have only one-dimensional \(H^1\), and Künneth
gives \(H^2(X,B)=\mathbf C\) with all other cohomology zero. Choose its
nonzero class in (E814.2). Equations (E814.4) and (E814.11) then give

\[
 (\operatorname{ext}^0,\operatorname{ext}^1,
   \operatorname{ext}^2,\operatorname{ext}^3,
   \operatorname{ext}^4)=(1,8,12,8,1),
 \quad\chi(F,F)=-2,
 \tag{E814.13}
\]

and full trace injectivity in degrees zero, one and two.

Duality preserves these properties. On \(X\times X\), Künneth and
gluability give

\[
 \operatorname{Ext}^2(F\boxtimes F^\vee,F\boxtimes F^\vee)
 =\bigoplus_{a+b=2}
   \operatorname{Ext}^a(F,F)\otimes
   \operatorname{Ext}^b(F^\vee,F^\vee).
\]

The full trace is the exterior tensor product on these summands. Their
targets have distinct Hochschild bidegrees \((-2,0),(-1,-1),(0,-2)\),
and on each summand both factors are injective. Hence the external
product is fully semiregular. It is simple and gluable by the same
Künneth argument. Its degree-two self-Ext dimension is \(12+8^2+12=88\).

\subsection{Scope relative to the split-eightfold
target}\label{SF-E814-scope-relative-to-the-split-eightfold-target}

The fourfold complex just constructed has the rational two-pure
character \(1-e^X\), with cap rank twelve. The required balanced-RM
character has cap rank eighteen. They are different characters. A fixed
Fourier--Mukai equivalence, isogeny pullback or rational algebraic
\(B\)-field shear preserves the relevant cap rank, so none turns this
fourfold complex into the required balanced-RM input. No claim is made
that its external product solves the eightfold problem either.

The role of this construction is precise: degree-two derived gluing is
genuinely available and can satisfy the full trace conditions, while the
coherent degree-one extensions excluded in SF-E813 cannot. A multivertex
construction with the balanced character still requires a new
calculation, including its finite coherent group action and full
degree-two trace.
